\emph{Preliminaries.} By Theorem \claimref{synthesis.signed-directions-and-the-integer-clifford-algebra} the algebra is free on the blades $e_{S}$, the images of the elements of the group $G$ of signed directions, and each blade is invertible with $e_{S}^{-1} = \pm e_{S}$. Inversion is an anti-automorphism of $G$ fixing the sign bit, so its linear extension $\dagger$ is an anti-automorphism of the algebra with $x^{\dagger\dagger} = x$. Since $e_{S}e_{T}^{\dagger} = \pm e_{S \triangle T}$ is a scalar only for $S = T$, where it is $1$, $\langle xy^{\dagger}\rangle_{0} = \sum_{S}x_{S}y_{S}$; this form is symmetric and $\langle xx^{\dagger}\rangle_{0} = \sum x_{S}^{2} \ge 0$. A self-adjoint blade satisfies $e_{S} = e_{S}^{\dagger} = e_{S}^{-1}$, so it has square $1$. Two blades commute or anticommute (Theorem \claimref{synthesis.blade-square-signs-f2-quadratic-form}(ii)).

\emph{(i).} $(A_{i}B_{j})^{\dagger} = B_{j}A_{i} = A_{i}B_{j}$, so $C^{\dagger} = C$. Write $C = A_{1}S + A_{2}D$ with $S = B_{1} + B_{2}$, $D = B_{1} - B_{2}$, which commute with the $A_{i}$. Then $C^{2} = S^{2} + D^{2} + A_{1}A_{2}SD + A_{2}A_{1}DS$, with $S^{2} + D^{2} = 4$, $SD = -[B_{1},B_{2}]$ and $DS = [B_{1},B_{2}]$; hence $C^{2} = 4 - [A_{1},A_{2}][B_{1},B_{2}]$. For a state and a scalar $t$, $0 \le \varphi((C - t)(C - t)^{\dagger}) = \varphi(C^{2}) - 2t\varphi(C) + t^{2}$; at $t = \varphi(C)$ this is $\varphi(C)^{2} \le \varphi(C^{2})$.

\emph{(ii).} \emph{The frame.} For each $(p,q) \in \mathcal{E}$ the self-adjoint blades are enumerated (the check does this for all fourteen): they are $1$, a set $F$ of pairwise anticommuting blades, and, in the five signatures named, the pseudoscalar $J$ (central, $J^{\dagger} = J$, $J^{2} = 1$) and the blades $Jh$, $h \in F$. Moreover the products $hk$ ($h \ne k$ in $F$, one for each unordered pair) are blades distinct from each other and from $1$ and the elements of $F$, also after multiplication by $J$. The sizes of $F$ are $0, 0$; $2, 2, 0$; $3, 2, 3, 0$; $5, 5, 5$; $5, 5$ in the order $(1,0), (0,1)$; $(2,0), (1,1), (0,2)$; $(3,0), (2,1), (1,2), (0,3)$; $(4,0), (1,3), (0,4)$; $(5,0), (1,4)$.

\emph{Blocks.} Without $J$ put $P = 1$. With $J$, the elements $P_{\pm} = \tfrac12(1 \pm J)$ are central self-adjoint idempotents with $P_{+} + P_{-} = 1$, $P_{+}P_{-} = 0$, and every statement below is made in the algebra $P_{\pm}\Cl(p,q)$ with unit $P_{\pm}$, where $Jh = \pm h$. In a block $P$, a self-adjoint element is $P(s + v)$ with $s$ a scalar and $v = \sum_{h \in F}v_{h}h$, uniquely; and since the $h$ anticommute and have square $1$, $v^{2} = Q(v) := \sum v_{h}^{2}$ and $vv' - v'v = 2\sum_{h < k}(v_{h}v'_{k} - v_{k}v'_{h})hk$.

\emph{Observables of a block.} $(s + v)^{2} = s^{2} + Q(v) + 2sv$. If this is $1$ in the block then $sv = 0$ (the blades of $F$ are independent of $1$ in the block), so either $v = 0$ and $s = \pm 1$, or $s = 0$ and $Q(v) = 1$. If two observables $v, v'$ with $s = s' = 0$ commute, then all $v_{h}v'_{k} - v_{k}v'_{h}$ vanish (the products $hk$ are independent in the block), so $v$ and $v'$ are proportional, and $Q(v) = Q(v') = 1$ gives $v' = \pm v$.

\emph{Conclusion.} Fix a block and suppose the components of $A_{1}, A_{2}$ in it do not commute. Then both have $s = 0$ and their vectors $v_{1}, v_{2}$ are not proportional. The component of $B_{j}$ commutes with both; if it had $s = 0$ its vector would be $\pm v_{1}$ and $\pm v_{2}$, which is impossible, so it is $\pm P$. Hence the components of $B_{1}, B_{2}$ commute. In every block one of the two commutators vanishes, so $[A_{1},A_{2}][B_{1},B_{2}] = 0$, $C^{2} = 4$ by (i), and $\varphi(C)^{2} \le \varphi(C^{2}) = 4$.

\emph{(iii).} \emph{The configurations.} For the thirteen signatures with $4 \le n \le 6$ outside $\mathcal{E}$ the check lists four self-adjoint blades with $A_{1}A_{2} = -A_{2}A_{1}$, $B_{1}B_{2} = -B_{2}B_{1}$ and $A_{i}B_{j} = B_{j}A_{i}$. For $n \ge 7$ choose six of the generators; the subalgebra they generate is $\Cl(p',q')$ with $p' + q' = 6$, which is not in $\mathcal{E}$, its blades are blades of $\Cl(p,q)$ with the same inverses, and a functional of the form $\langle\rho\,\cdot\,\rangle_{0}$ with $\rho$ in the subalgebra is defined on all of $\Cl(p,q)$; so the configuration and the states below carry over. Every signature with $n \ge 4$ outside the five listed is covered.

\emph{The square.} $[A_{1},A_{2}] = 2A_{1}A_{2}$ and $[B_{1},B_{2}] = 2B_{1}B_{2}$, so by (i) $C^{2} = 4 + 4K$ with $K = -A_{1}A_{2}B_{1}B_{2}$. $K$ is a blade up to sign; $K^{\dagger} = -B_{2}B_{1}A_{2}A_{1} = -A_{1}A_{2}B_{1}B_{2} = K$ (two sign changes), so $K^{2} = 1$; $A_{1}$ anticommutes with $A_{1}A_{2}$ and commutes with $B_{1}B_{2}$, and likewise for the others, so each $A_{i}B_{j}$ commutes with $K$ and $KC = CK$. $K$ is not a scalar multiple of $1$: otherwise $A_{1}A_{2} = \pm B_{1}B_{2}$ would commute with $A_{1}$. Since the $A$'s commute with the $B$'s, $KA_{1}B_{1} = -(A_{1}A_{2}A_{1})(B_{1}B_{2}B_{1}) = -(-A_{2})(-B_{2}) = -A_{2}B_{2}$, and in the same way $KA_{2}B_{2} = -A_{1}B_{1}$, $KA_{1}B_{2} = A_{2}B_{1}$, $KA_{2}B_{1} = A_{1}B_{2}$; so $KC = C$, and $C^{3} = 4C + 4KC = 8C$. Each term $A_{i}B_{j}$ is a blade that is not a scalar ($A_{i} = \pm B_{j}$ is excluded because $B_{j}$ commutes with both $A$'s), so $\langle C\rangle_{0} = \langle KC\rangle_{0} = 0$, and $\langle K\rangle_{0} = 0$.

\emph{The functionals.} $\rho_{r} = (1 + K)(1 + rC)$ is self-adjoint. $\varphi_{r}(1) = \langle 1 + K + rC + rKC\rangle_{0} = 1$; $\varphi_{r}(C) = \langle(1 + K)C\rangle_{0} + r\langle(1 + K)C^{2}\rangle_{0} = 0 + r\langle(1 + K)(4 + 4K)\rangle_{0} = 8r$; $\varphi_{r}(C^{2}) = \langle(1 + K)(4 + 4K)\rangle_{0} + r\langle(1+K)C^{3}\rangle_{0} = 8 + 8r\langle(1 + K)C\rangle_{0} = 8$. Positivity: put $P = \tfrac12(1 + K)$, a self-adjoint idempotent commuting with $C$, and $y = Px$. Since $\langle ab\rangle_{0} = \langle ba\rangle_{0}$, $\tfrac12\varphi_{r}(xx^{\dagger}) = \langle x^{\dagger}P(1 + rC)Px\rangle_{0} = \langle y^{\dagger}y\rangle_{0} + r\langle y^{\dagger}Cy\rangle_{0}$. By the Cauchy-Schwarz inequality for the positive form $\langle a^{\dagger}b\rangle_{0}$ (an inequality between squares, valid in any ordered field), $\langle y^{\dagger}Cy\rangle_{0}^{2} \le \langle y^{\dagger}y\rangle_{0}\langle y^{\dagger}C^{2}y\rangle_{0} = 8\langle y^{\dagger}y\rangle_{0}^{2}$, because $C^{2}P = 8P$. So $(r\langle y^{\dagger}Cy\rangle_{0})^{2} \le 8r^{2}\langle y^{\dagger}y\rangle_{0}^{2} \le \langle y^{\dagger}y\rangle_{0}^{2}$ when $8r^{2} \le 1$, and $\varphi_{r}(xx^{\dagger}) \ge 0$. Conversely if $\varphi_{r}$ is a state then by (i) $64r^{2} = \varphi_{r}(C)^{2} \le \varphi_{r}(C^{2}) = 8$. The set of rational $r$ with $8r^{2} \le 1$ has supremum $\sqrt{2}/4$, which is not rational, so the values $8r$ have supremum $2\sqrt{2}$, not attained; a field containing $\sqrt{2}$ contains $r = \sqrt{2}/4$. That no state exceeds $2\sqrt{2}$ in any signature is Tsirelson's bound and is not needed for this result.

\emph{(iv).} By Theorem \claimref{transport.classification-real-clifford-algebras} $\Cl(p,q)$ is $\Mat_{k}(D)$ or a sum of two copies, $D$ one of $\R$, $\C$, $\HH$ according to $p - q \bmod 8$, with $k^{2}\dim D$ (times $2$ for a sum) equal to $2^{n}$. Computing $k$ from the table for each signature gives $k \le 2$ exactly on $\mathcal{E}$ for $n \le 6$ (the check does this up to $n = 8$), and $k \ge 4$ for $n \ge 7$ since $k^{2} \ge 2^{n}/8$. The named cases: $\Cl(3,0) = \Mat_{2}(\C)$, $\Cl(1,3) = \Mat_{2}(\HH)$, $\Cl(3,1) = \Cl(2,2) = \Mat_{4}(\R)$.
